\documentclass[10pt,twocolumn]{article}
\usepackage[a4paper,margin=15mm]{geometry}
\usepackage{amsmath,amssymb}
\usepackage{mathrsfs}
\usepackage{amsthm}
\theoremstyle{definition}
\newtheorem{theorem}{Theorem}
\newtheorem{proposition}{Proposition}
\newtheorem{corollary}{Corollary}
\newtheorem{remark}{Remark}
\newtheorem{definition}{Definition}
\usepackage{graphicx}
\usepackage{booktabs,array,calc}
\usepackage[font=small,labelfont=bf]{caption}
\usepackage[colorlinks=true]{hyperref}
\usepackage{etoolbox}
\providecommand{\tightlist}{\setlength{\itemsep}{0pt}\setlength{\parskip}{0pt}}
\setcounter{secnumdepth}{-1}
\emergencystretch=3em
\allowdisplaybreaks
\AtBeginDocument{%
  \BeforeBeginEnvironment{equation}{\small}%
  \BeforeBeginEnvironment{equation*}{\small}%
  \BeforeBeginEnvironment{align}{\small}%
  \BeforeBeginEnvironment{align*}{\small}%
  \BeforeBeginEnvironment{gather}{\small}%
  \BeforeBeginEnvironment{gather*}{\small}%
}

\begin{document}
\title{Obstruction Certificates as Dual Objects:\\ One Template, Four Instantiations}
\author{Amin Abaee\\[0.35em]
{\small Independent Researcher}\\[0.55em]
{\small\href{https://orcid.org/0000-0002-0019-1842}{ORCID: 0000-0002-0019-1842}}\\[0.3em]
{\small\href{mailto:amin\_abaee@ut.ac.ir}{amin\_abaee@ut.ac.ir}}}
\date{September 23, 2026}
\maketitle

\begin{abstract}
The obstruction calculus issues four families of nonviability certificates
--- common-action, timing, fibre-certification, and certainty-equivalence
--- that were derived separately. This paper, the D6 item of the
generalization catalogue, unifies their computational reading: each
certificate is a \emph{dual feasibility object} for the primal system the
corresponding obstruction declares infeasible. The common-action
certificate is a Farkas multiplier vector; the timing certificate is a
telescoping (additive) or substitution (multiplicative) identity on the
decline rows; the fibre certificate is a contradicting-witness pair inside
one observation cell; the certainty-equivalence certificate is an affine
drift lower bound. On the programme's working instances all four dual
objects are computed and validated in exact rational arithmetic, together
with the template's scope statement: duals certify nonviability only,
viability being certified primitively by witness actions --- and at the
corner state where no witness exists, primal and dual agree in declaring
nonviability (6/6 exact check families).
\end{abstract}

\section{Introduction and relation to the programme}

The completed items derive their certificates mechanism by mechanism. The
source catalogue proposes the applied-mathematics reading adopted here:
\emph{the obstruction calculus produces primal-dual certificates of
nonviability under information constraints}. The paper makes that reading
exact on the programme's own audited instances --- no new mechanism is
introduced, and no mechanism from the completed items is modified; the
contribution is the identification, the exact witnesses, and the scope
statement.

\section{The template}

Each obstruction family declares a primal feasibility problem infeasible:
\begin{itemize}\tightlist
\item \textbf{Common action:} the stacked safe-action system
$Au\le b$ (one instrument must serve several states' floors);
\item \textbf{Timing:} the decline-chain program $z_t\ge1$,
$z_{t+1}=z_t-\varepsilon$ (or its multiplicative contraction
$q_{t+1}=\tfrac9{10}q_t$);
\item \textbf{Fibre certification:} a single exact label for all states of
one observation cell;
\item \textbf{Certainty equivalence:} a law in the CE class whose
corrected drift stays below the viability threshold.
\end{itemize}
The template: an obstruction certificate \emph{is} a dual object --- a
multiplier vector, a telescoping identity, a witness pair, an affine lower
bound --- each checkable in exact rational arithmetic against its own
primal.

\begin{proposition}[Common action]
$Au\le b$ is infeasible iff a rational $\lambda\ge0$ exists with
$\lambda^{\top}A=0$, $\lambda^{\top}b<0$; the multiplier vector is the
obstruction certificate (Farkas' lemma).
\end{proposition}

\begin{proposition}[Timing]
\label{prop:timing}
For the additive chain, summing the $K$ decline rows with all-ones
multipliers yields the identity $z_0-z_K=K\varepsilon$, whence survival
for $K$ steps holds iff $z_0\ge1+K\varepsilon$; for the contraction
chain, substitution yields $q_K=q_0(\tfrac9{10})^K$, whence survival iff
$q_0\ge(\tfrac{10}9)^K$ --- the level sets of the successor paper's
exit-time value. The dual objects are the identity rows themselves.
\end{proposition}
\begin{proof}
Telescoping sums of equalities are equalities; the multiplic case is the
same statement with multiplication. Both directions are exact, so the
threshold identities hold with equality at the boundary.
\end{proof}

\begin{proposition}[Fibre and CE]
If one observation cell carries two states whose exact labels differ, the
cell admits no exact certifier, and the dual object is the contradicting
pair. If every law in a policy class has drift bounded below by an affine
function exceeding the viability threshold on the audited grid, the class
is empty as a correspondence with safety, and the affine bound is the dual
object.
\end{proposition}
\begin{proof}
Immediate from the definitions; the content is the exact instances below.
\end{proof}

\section{Verified instances (6/6 exact check families)}

\paragraph{E1.} For the two-floor stacked system $u\le\tfrac25\wedge
u\ge\tfrac35$: $\lambda=(1,1)$, $\lambda^{\top}A=0$,
$\lambda^{\top}b=-\tfrac15<0$.

\paragraph{E2.} Additive thresholds $z_0\ge1+K/10$ ($K\le4$) and
multiplicative thresholds $q_0\ge(\tfrac{10}9)^K$ ($K\le3$) verified at
exact powers and $\tfrac1{100}$-neighbors --- equality at the boundary in
both cases, matching Proposition~\ref{prop:timing}.

\paragraph{E3.} The crossing fibre of the two-floor index instance carries
$s=\tfrac1{10}$ (label low) and $s=\tfrac25$ (label high) in one cell
$\mathcal O=1$: no exact certifier exists; the dual object is the pair.

\paragraph{E4.} On the CE trap with reading $g(s)=s^2$ and correction
$\tfrac1{10}$: the drift is exactly $g(s+\tfrac1{10})-g(s)=\tfrac{s}5+
\tfrac1{100}$ on the $101$-point grid, and $\tfrac{s}5+\tfrac1{100}\ge
\tfrac{21}{100}$ iff $s\ge1$ --- the affine dual bound is exact.

\paragraph{E5 (template closure).} Every dual object re-validates its
primal infeasibility by direct evaluation; the feasible contrast system
($u\le\tfrac25\wedge u\ge\tfrac15$) admits the witness $u=\tfrac15$ and no
certificate: primal and dual agree on both sides of every verdict.

\paragraph{E6 (scope).} On the two-floor two-patch system, safe states off
the corner $(1,1)$ have primal witness actions; at the corner both primal
(one-step exit under every instrument) and dual (its stacked system is
trivially infeasible) certify nonviability. Duals certify nonviability
only; viability is never dual-certified.

\section{What is not claimed}

The template is an identification, not a new algorithm: no solver is
developed, no complexity claim is made, and the barrier-style/SOS dual
objects mentioned in the source catalogue for nonlinear systems are cited
as outlook, not computed (the programme's instances are discrete and
rational). Sum-of-squares methods, semidefinite duals, and continuous-state
duality are out of scope. The four instantiations share a lemma (Farkas)
only at the level of template; no deeper converse is claimed --- in
particular the paper does not claim that every nonviability proof in the
programme must dualize.


\section*{Declarations}
\textbf{Funding.} None.\quad \textbf{Competing interests.} None.\quad
\textbf{Data availability.} No external data were used.\quad
\textbf{Code availability.} The verification script
\texttt{certificate\_duality\_v1\_verify.py} is publicly available at
\url{https://github.com/MIKEAA2020/general-sustainability} (folder
\texttt{arena agent 1/paper rewrites/latex}); it reproduces all six check
families verbatim.\quad
\textbf{AI declaration.} GLM (Z.ai), Qwen (Alibaba Cloud) and DeepSeek AI
assisted with drafting and iterative review.

\end{document}
